\section{Conclusion}
\label{sec:conclusion}

This paper develops a comprehensive econophysics pipeline for analysing the dependency structure of Brazilian equities traded on B3 from 1998 to 2025. The main findings can be summarised along four dimensions.

First, the stylized-fact and correlation analysis confirms that Brazilian equities display the empirical signatures expected from complex financial systems: fat-tailed return distributions, persistent volatility clustering and a moderate but positive average cross-correlation ($\bar{\rho} \approx 0.24$). Within-sector correlations are statistically significantly higher than between-sector correlations, providing early evidence that sectoral organisation is embedded in the empirical dependency matrix.

Second, Random Matrix Theory separates the correlation matrix into interpretable components. Five eigenvalues exceed the Mar\v{c}enko--Pastur upper bound, with the largest eigenvalue accounting for approximately $37.3\%$ of the total variance. The Market Mode captures broad co-movement, while the remaining significant eigenvectors reveal commodity, financial and utility-related group structures. The spectral reconstruction is numerically stable, and the group-mode matrix retains meaningful localized dependency blocks after the dominant market component is removed.

Third, financial network analysis through MST and PMFG representations shows that RMT filtering changes the topology of asset dependencies. The group-mode PMFG has lower mean edge correlation but higher clustering than the original PMFG, indicating that market-mode removal reveals denser localized communities. Hub rankings shift from market-wide connectors (GGBR4, BBDC4) to localized sectoral bridges (GUAR3). The PMFG retains 168 edges with 166 triangles and 55 four-cliques, providing richer local geometry than the MST and enabling more detailed mesoscopic analysis.

Fourth, the volatility forecasting experiment provides evidence that network topology can complement classical volatility predictors. GARCH(1,1) and HAR-RV remain strong benchmarks, especially at short horizons, consistent with the high persistence ($\alpha + \beta > 0.98$) observed in all three target assets. Machine-learning models with PMFG-derived features approach these baselines and, at the 20-day horizon, Ridge Regression slightly outperforms GARCH in QLIKE. Feature importance analysis shows that while classical autoregressive features dominate, PMFG centralities contribute nonzero predictive information. The appropriate conclusion is incremental: network features do not replace classical volatility dynamics, but they add structural information from the cross-sectional geometry of the market.

Future work should extend this framework in several directions: constructing rolling RMT and network features for fully real-time forecasting, expanding the target asset set beyond the three demonstration equities, comparing against multivariate volatility models such as DCC-GARCH, and evaluating temporal neural network and graph neural network architectures that can jointly exploit the spatial structure of financial networks and the temporal dynamics of volatility. The broader goal is to connect interpretable complex-systems descriptions of financial dependency with practical risk forecasting in emerging markets.
